\chapter{Introduction}
\label{ch:intro}

\section{Background}
\label{sec:background}

Visuomotor policies such as Action Chunking Transformers (ACT) and Diffusion
Policy have made behaviour cloning a practical route to robot manipulation, but
their performance scales with the quantity and quality of demonstration data.
The dominant way of acquiring that data is teleoperation: a human operator
drives the robot through the task while joint states and camera streams are
recorded. Teleoperation yields action labels that are, by construction, exactly
executable on the target embodiment, which is its principal advantage. Its
disadvantages are equally structural. Each episode consumes an operator's time
in real time, the required hardware---a VR controller, an exoskeleton or a
haptic device---is expensive and embodiment-specific, and the resulting motion
is shaped as much by the latency and ergonomics of the interface as by the task
itself.

Human demonstrations recorded passively, with the demonstrator using their own
hands and no interface at all, invert this trade-off. Data collection becomes
cheap, fast and natural; a demonstrator can record dozens of episodes in the
time a teleoperator records one, and the recorded motion reflects how the task
is actually performed rather than how it can be performed through a control
interface. What is lost is precisely the property teleoperation guarantees:
there are no action labels, and the kinematics of a human arm and hand do not
correspond to the kinematics of a robot manipulator. Bridging that gap is the
retargeting problem, and it is the subject of this thesis.

\section{Problem Statement}
\label{sec:problem}

Passive capture makes demonstration data cheap but produces observations, not
actions. Converting an observed human hand trajectory into a robot joint
trajectory that (i) reproduces the task-relevant motion, (ii) is smooth enough
for a policy to imitate, and (iii) is physically executable on the target
manipulator, is an underdetermined problem contaminated by perception noise.

Two failure modes dominate. The first is noise amplification through the
kinematic chain. Monocular or neural depth estimates of hand keypoints are
noisy at the centimetre scale; inverse kinematics is nonlinear, so a bounded
task-space perturbation may map to an unbounded joint-space excursion near
singularities. The widespread response---low-pass filtering the observed
trajectory before solving inverse kinematics, then filtering the joint
trajectory afterwards---does not solve this. Task-space pre-filtering is
applied in the wrong space to guarantee joint-space smoothness, and it removes
genuine high-frequency task content along with the noise. This is the
garbage-in-garbage-out pattern that motivates the present work.

The second failure mode is representational. A trajectory recorded as an
absolute sequence of positions in a camera or world frame is tied to the exact
spatial layout of the recording session. If the object is moved, the trajectory
is wrong; if the embodiment changes, it is wrong in a different way. What
generalises is the relation between the hand and the object being manipulated,
not the absolute path of the hand through space.

\section{Research Gap}
\label{sec:gap}

Both failure modes have known remedies, and the literature has adopted them---
but in disjoint communities.

Optimisation-based retargeting, in which data fidelity, temporal smoothness and
joint limits are posed as a single constrained optimisation over the robot
configuration, is standard in \emph{online teleoperation}. AnyTeleop and the
associated \texttt{dex-retargeting} library established the canonical form;
Open-TeleVision, Bunny-VisionPro and ACE build directly on it; DexFlow extends
the smoothness penalty to second order. In every case, however, a human remains
in the control loop, the optimisation runs at interactive rates, and no dataset
is produced.

Conversely, \emph{offline} pipelines that generate datasets from passive human
video adopt object-centric representations but not the joint-optimisation
formulation. ORION plans over sparse object keyframes and does not perform
per-frame hand retargeting at all. OKAMI reconstructs a parametric body-and-hand
model and retargets in two decoupled stages, warping a reference trajectory to
a new object position and then solving inverse kinematics---structurally the
same sequential pattern criticised in Section~\ref{sec:problem}.

A targeted prior-art review conducted for this work (Chapter~\ref{ch:lit},
Section~\ref{sec:positioning}) found no published system combining passive
markerless capture, a metrically calibrated multi-camera RGB-D rig, an
object-centric dense trajectory representation, single-cost-functional
retargeting, imitation-learning-ready dataset export, and physical replay used
as a pre-export validation step. The nearest systems each satisfy some criteria
and miss others. The gap addressed here is therefore not the invention of a new
retargeting operator, but the transfer of an established formulation into a
regime where it has not been applied, together with the sensing and validation
infrastructure that regime requires.

\section{Research Questions and Hypotheses}
\label{sec:rq}

The work is guided by one primary research question and three sub-questions.

\vspace{2mm}
\noindent\textbf{RQ.} \emph{How can human manipulation demonstrations captured
passively by a calibrated multi-camera RGB-D rig be retargeted onto a robot
manipulator so as to yield datasets suitable for training visuomotor imitation
policies?}

\begin{enumerate}[label=\textbf{RQ\arabic*.}, leftmargin=2.2em]
  \item How does formulating retargeting as a single weighted cost functional,
        combining a confidence-weighted data term with velocity and
        acceleration penalties, affect the fidelity and smoothness of the
        resulting joint trajectories compared with a sequential
        smooth--IK--smooth pipeline?
  \item What accuracy can be achieved for wrist pose estimation by
        confidence-weighted fusion of multiple calibrated RGB-D views anchored
        to a common ChArUco workspace frame, and how does the residual error
        decompose across its independent sources?
  \item To what extent does replaying synthesised trajectories on the physical
        manipulator, prior to dataset export, improve the executability of the
        exported dataset?
\end{enumerate}

\noindent The corresponding hypotheses are stated below. Following the
convention for experimental computer-science hypotheses, each identifies an
independent variable, a dependent variable, a population and a relation.

\begin{description}[leftmargin=3.2em, style=nextline]
  \item[\textnormal{\textbf{H1 (behaviour).}}]
  Retargeting posed as a single weighted cost functional produces joint
  trajectories with lower task-space tracking error and higher smoothness than
  a sequential smooth--IK--smooth pipeline, on the same recorded demonstrations.
  \emph{Independent variable:} retargeting formulation (joint vs.\ sequential).
  \emph{Dependent variables:} wrist-pose tracking RMSE; spectral arc length of
  the joint trajectory. \emph{Population:} recorded demonstrations of the
  constrained task class defined in Section~\ref{sec:tasks}.
  \emph{Relation:} the joint formulation dominates on both measures.

  \item[\textnormal{\textbf{H2 (behaviour).}}]
  Weighting the data term by per-frame perception confidence reduces tracking
  error relative to uniform weighting, and the reduction is larger on frames
  where hand landmarks are partially occluded.
  \emph{Independent variable:} presence of confidence weighting.
  \emph{Dependent variable:} wrist-pose tracking RMSE, stratified by occlusion.
  \emph{Population:} as above. \emph{Relation:} monotone reduction, amplified
  under occlusion.

  \item[\textnormal{\textbf{H3 (dependability).}}]
  Screening trajectories by physical replay before export increases the
  proportion of exported episodes that execute successfully on the manipulator.
  \emph{Independent variable:} presence of the replay screening stage.
  \emph{Dependent variable:} executable fraction of the exported dataset.
  \emph{Population:} as above. \emph{Relation:} positive.
\end{description}

The questions were checked against the FINER criteria. They are
\emph{feasible}, in that the hardware, solver and evaluation protocol are all
available and the task class is deliberately constrained; \emph{interesting}
and \emph{relevant} to the robot-learning community, because demonstration cost
is a recognised bottleneck; \emph{novel} in the specific sense established in
Section~\ref{sec:gap}, namely a regime transfer rather than a new operator; and
\emph{ethical}, as the study involves no human subjects beyond the author
performing demonstrations and no personally identifying data is retained.

\section{Contributions}
\label{sec:contributions}

This thesis makes the following contributions.

\begin{enumerate}[leftmargin=2em]
  \item A retargeting formulation for the \emph{offline} setting that poses
        data fidelity, velocity smoothness and acceleration smoothness as one
        constrained optimisation over the robot configuration, with the data
        term weighted by per-frame perception confidence and made robust by a
        Huber loss (Section~\ref{sec:retarget}).
  \item An object-centric dense trajectory representation that stores each
        demonstration both relative to the manipulated object and relative to a
        fixed ChArUco workspace anchor, permitting downstream consumers to
        select the representation that generalises better
        (Section~\ref{sec:objcentric}).
  \item A multi-camera RGB-D capture and calibration procedure with
        confidence-weighted depth fusion and an explicit treatment of
        cross-device timestamp alignment, together with an error budget that
        separates the three independent contributions to residual pose error
        (Sections~\ref{sec:calib}--\ref{sec:fusion}, Section~\ref{sec:errbudget}).
  \item A physical-replay validation stage that filters kinematically
        infeasible trajectories before dataset export, and an empirical
        characterisation of what that stage rejects
        (Sections~\ref{sec:replay}, \ref{sec:res-replay}).
  \item An open implementation exporting LeRobot-compatible datasets, evaluated
        end-to-end by training ACT and Diffusion Policy on the generated data
        (Chapter~\ref{ch:impl}).
\end{enumerate}

\section{Thesis Structure}
\label{sec:structure}

Chapter~\ref{ch:lit} reviews prior work on hand representation, retargeting
formulations, task representation, multi-camera RGB-D calibration and dataset
generation, and positions \viki\ against the nearest published systems.
Chapter~\ref{ch:method} describes the design of the system and the derivation
of the retargeting cost functional. Chapter~\ref{ch:impl} reports the
implementation, the evaluation protocol and the experimental results.
Chapter~\ref{ch:disc} interprets those results, decomposes the error budget and
states the limitations. Chapter~\ref{ch:concl} concludes and outlines future
work.
